\section{Related Work}
\label{sec:related_work}

Most self-supervised methods for images learn representations by making the
embeddings of two augmented views of the same image agree, while avoiding the
trivial solution in which every image maps to the same embedding (complete
collapse). Contrastive methods such as SimCLR~\cite{chen2020simclr} and
MoCo-v3~\cite{chen2021mocov3} prevent collapse by also pushing apart
embeddings of different images (negative pairs); self-distillation methods
such as BYOL~\cite{grill2020byol} and DINO~\cite{caron2021dino} prevent it
with asymmetry between two networks (momentum encoders, predictor heads);
masked autoencoders~\cite{he2022mae} avoid view agreement altogether by
reconstructing masked input patches. Barlow Twins~\cite{zbontar2021barlow}
instead pushes the cross-correlation matrix between the embeddings of the two
views toward the identity matrix: the diagonal terms enforce view invariance,
while the off-diagonal terms decorrelate embedding dimensions, so the
objective itself penalizes collapse, with neither negative pairs nor network
asymmetry. We adopt Barlow Twins unchanged --- reference loss implementation
and published augmentations --- and characterize how it behaves as
pre-training data shrinks.

How much pre-training data self-supervised learning needs is a well-studied
question. Contrastive learning has been measured across data quantity,
domain, and quality: the benefit of pre-training data beyond 500k images is
modest~\cite{cole2022when}. Closer to our scale, denoising autoencoders
(BEiT-style masked prediction) have been reported more robust to the size and
type of the pre-training set than joint-embedding methods --- represented
there by DINO --- when evaluated by fine-tuning; under linear probing, the
same study reports denoising autoencoders falling behind joint-embedding
methods~\cite{elnouby2021largescale}. At the opposite end of the scale,
joint-embedding training on a large enough dataset needs almost no
augmentations beyond cropping: hand-crafted augmentations act chiefly by
artificially enlarging the dataset rather than by enforcing useful
invariances, so their benefit concentrates where data is
limited~\cite{moutakanni2024you} --- though the smallest scale tested there
is 1.3M images, far above ours. These studies chart
the small-data regime of self-supervised learning across method families, but
none of them measures Barlow Twins itself in that regime --- so we
characterize the data efficiency of Barlow Twins directly, on a small ViT,
from 1k to 100k pre-training images, under a linear-probe evaluation.

Comparing models pre-trained on different set sizes requires a training
schedule that does not tie model quality to a pre-committed stopping point.
Cosine learning-rate schedules~\cite{loshchilov2017sgdr} anneal toward a
fixed horizon, and a checkpoint read out well before that horizon attains
higher loss than a model whose schedule was matched to that
duration~\cite{hoffmann2022chinchilla}. Warmup-stable-decay (WSD) schedules
--- linear warmup, a long constant-rate phase, and a short decay --- match
the final loss of a well-tuned cosine schedule without committing to a
horizon~\cite{hagele2024scaling,hu2024minicpm}, with a $(1-\sqrt{t})$ decay
shape the best performer in the cooldown-shape ablation of H\"agele et
al.~\cite{hagele2024scaling}. Schedules of this family also have vision
precedent: vision transformers have been scaled with constant or
reciprocal-square-root phases followed by a cooldown, coming close to a
horizon-committed schedule while letting one run serve all training
durations~\cite{zhai2022scalingvit}, and horizon-free schedules are effective
for continual pre-training, including masked-autoencoder vision
models~\cite{singh2025beyond}. So we train every pre-training set size
with a standard WSD-style schedule, assembled from these published components
(Section~\ref{sec:method}); the schedule is not a contribution of this work.

Monitoring self-supervised pre-training requires a validation signal without
labels for the pretext task. Classifying held-in-domain images with a
$k$-nearest-neighbour ($k$NN) probe on frozen features is the standard such
monitor, used by DINO with $k{=}20$~\cite{caron2021dino}. The effective rank
of the representation has been shown to track downstream accuracy, enabling
fully label-free model selection~\cite{garrido2023rankme}; we log this signal
but do not select with it. Selecting the best value of a repeatedly evaluated
validation signal inflates the selected estimate --- over-fitting in model
selection~\cite{cawley2010overfitting}. So we select checkpoints by an
in-domain $k$NN probe in the DINO style and smooth the validation curve
before taking its maximum, to reduce selection bias.

A representation can degrade without collapsing completely: in dimensional
collapse, the embedding variance concentrates in a low-dimensional
subspace~\cite{jing2022dimensional}, which the effective rank of the feature
covariance spectrum quantifies~\cite{roy2007effective}. Augmentations are not
a free remedy in this regime: the invariances a joint-embedding method learns
are the invariances its augmentations define~\cite{xiao2021what}, and views
that share too little information destroy task-relevant
signal~\cite{tian2020infomin}; for supervised ViTs, augmentation and
regularization can substitute for an order of magnitude of training
data~\cite{steiner2022train}, underlining how strongly augmentation choices
act as implicit data. So in the small-data regime we keep the published
augmentation recipe fixed, track effective rank as a logged collapse
diagnostic, and instead vary a capacity knob that is specific to Barlow Twins
--- the projector width.
